\documentclass[12pt,a4paper]{article}
\usepackage[utf8]{inputenc}
\usepackage[T1]{fontenc}
\usepackage[spanish]{babel}
\usepackage{geometry}
\usepackage{hyperref}
\usepackage{graphicx}
\usepackage{listings}
\usepackage{xcolor}
\usepackage{booktabs}
\usepackage{tabularx}
\usepackage{caption}
\usepackage{float}
\usepackage{parskip}
\usepackage{enumitem}
\usepackage{titlesec}

\geometry{margin=2.5cm}

% Cambiar "Cuadro" (babel español) por "Tabla"
\addto\captionsspanish{\renewcommand{\tablename}{Tabla}}

% Cambiar "Listing" por "Listado"
\renewcommand{\lstlistingname}{Listado}

\hypersetup{
  colorlinks=true,
  linkcolor=blue!60!black,
  urlcolor=blue!60!black,
  citecolor=blue!60!black
}

\definecolor{codegray}{rgb}{0.95,0.95,0.95}
\definecolor{codecomment}{rgb}{0.4,0.4,0.4}
\definecolor{codekeyword}{rgb}{0.0,0.3,0.7}
\definecolor{codestring}{rgb}{0.2,0.5,0.2}
\definecolor{codeannotation}{rgb}{0.6,0.3,0.0}

\lstset{
  language=Java,
  backgroundcolor=\color{codegray},
  basicstyle=\ttfamily\small,
  keywordstyle=\color{codekeyword}\bfseries,
  commentstyle=\color{codecomment}\itshape,
  stringstyle=\color{codestring},
  numbers=left,
  numberstyle=\tiny\color{gray},
  numbersep=8pt,
  breaklines=true,
  showstringspaces=false,
  frame=single,
  framerule=0pt,
  xleftmargin=1em,
  xrightmargin=0.5em,
  captionpos=b,
  extendedchars=true,
  literate=
    {á}{{\'a}}1 {é}{{\'e}}1 {í}{{\'\i}}1 {ó}{{\'o}}1 {ú}{{\'u}}1
    {Á}{{\'A}}1 {É}{{\'E}}1 {Í}{{\'I}}1 {Ó}{{\'O}}1 {Ú}{{\'U}}1
    {ñ}{{\~n}}1 {Ñ}{{\~N}}1 {ü}{{\"u}}1 {Ü}{{\"U}}1
    {@}{{{\color{codeannotation}@}}}1
}

\title{\textbf{Herramientas de testing en OpenMarkov}\\[0.5em]
\large Informe de referencia}
\author{Manuel Arias Calleja --- CISIAD, UNED}
\date{Abril 2026}

\begin{document}

\maketitle
\tableofcontents
\newpage

% ─────────────────────────────────────────────────────────────────────────────
\section{Introducción}
\label{sec:intro}

Este documento describe las herramientas de testing disponibles en OpenMarkov,
tanto las que ya se usaban antes del esfuerzo actual de cobertura como las que
se han incorporado recientemente. El objetivo es servir de referencia rápida
para cualquier desarrollador que necesite escribir o interpretar tests en el
proyecto.

La Tabla~\ref{tab:resumen} resume el conjunto completo de herramientas.

\begin{table}[H]
\centering
\begin{tabularx}{\textwidth}{llcX}
\toprule
\textbf{Herramienta} & \textbf{Versión} & \textbf{Origen} & \textbf{Propósito} \\
\midrule
JUnit 5 (JUnit Platform + Jupiter) & 5.13.0 & Preexistente & Plataforma y anotaciones de tests \\
AssertJ Core & 4.0.0-M1 & Preexistente & Aserciones fluidas \\
AssertJ Swing & 3.17.1 & Preexistente & Tests de interfaz gráfica Swing \\
jqwik & 1.9.2 & Nuevo & Tests basados en propiedades \\
PITest (pitest-maven) & 1.17.3 & Nuevo & Testing de mutaciones \\
pitest-junit5-plugin & 1.2.0 & Nuevo & Integración PITest con JUnit 5 \\
JaCoCo & 0.8.11 & Preexistente & Cobertura de código \\
Maven Surefire & 3.5.2 & Preexistente & Ejecución de tests en Maven \\
\bottomrule
\end{tabularx}
\caption{Herramientas de testing en OpenMarkov}
\label{tab:resumen}
\end{table}

% ─────────────────────────────────────────────────────────────────────────────
\section{Herramientas preexistentes}
\label{sec:preexistentes}

% ─────────────────────────────────────────────────────────────────────────────
\subsection{JUnit 5}
\label{sec:junit5}

JUnit 5 es la plataforma de ejecución de tests del proyecto. Está gestionado
mediante un BOM (\emph{Bill of Materials}) en el POM padre, lo que garantiza
versiones coherentes en todos los módulos.

La anotación fundamental es \texttt{@Test}, que marca un método como caso de
prueba. JUnit se encarga de descubrir estos métodos, ejecutarlos y reportar
los resultados.

\begin{lstlisting}[caption={Test básico con JUnit 5}, label={lst:junit-test}]
import org.junit.jupiter.api.Test;
import static org.assertj.core.api.Assertions.*;

class VariableTest {

    @Test
    void numStatesMatchesConstructorArgument() {
        Variable v = new Variable("Rain", 3);
        assertThat(v.getNumStates()).isEqualTo(3);
    }
}
\end{lstlisting}

Otras anotaciones de JUnit 5 usadas en el proyecto:

\begin{description}[leftmargin=2em, style=nextline]
  \item[\texttt{@BeforeEach}] Método ejecutado antes de \emph{cada} test.
    Se usa para inicializar el estado compartido (por ejemplo, crear una red
    \texttt{ProbNet} fresca).
  \item[\texttt{@AfterEach}] Método ejecutado después de cada test.
    Se usa para liberar recursos.
  \item[\texttt{@Disabled}] Desactiva un test temporalmente sin eliminarlo.
    Surefire lo contabiliza como \emph{skipped}.
  \item[\texttt{@Nested}] Agrupa tests relacionados en clases internas,
    creando una jerarquía legible en los informes.
\end{description}

% ─────────────────────────────────────────────────────────────────────────────
\subsection{AssertJ Core}
\label{sec:assertj}

AssertJ proporciona una API fluida de aserciones que sustituye a los
\texttt{assertEquals}/\texttt{assertTrue} clásicos de JUnit. Su principal
ventaja es la legibilidad y los mensajes de error detallados.

\begin{lstlisting}[caption={Aserciones fluidas con AssertJ}, label={lst:assertj}]
// Comparación simple
assertThat(node.getParents()).contains(parentNode);

// Comparación numérica con tolerancia
assertThat(sum).isCloseTo(1.0, within(1e-10));

// Verificar excepciones
assertThatThrownBy(edit::executeEdit)
    .isInstanceOf(ConstraintViolatedException.class);

// Verificar identidad de objeto (misma referencia)
assertThat(restored).isSameAs(original);

// Verificar contenido exacto de una colección
assertThat(copy.getValues()).containsExactly(original.getValues());
\end{lstlisting}

Los métodos más usados en el proyecto son:

\begin{description}[leftmargin=2em, style=nextline]
  \item[\texttt{isEqualTo(expected)}] Igualdad por \texttt{equals()}.
  \item[\texttt{isSameAs(expected)}] Identidad de referencia (\texttt{==}).
  \item[\texttt{isNull()} / \texttt{isNotNull()}] Verificar nulidad.
  \item[\texttt{contains(element)}] La colección contiene el elemento.
  \item[\texttt{containsExactly(elements...)}] Contenido y orden exactos.
  \item[\texttt{doesNotContain(element)}] Ausencia en una colección.
  \item[\texttt{isCloseTo(value, within(delta))}] Comparación numérica con
    tolerancia, imprescindible para valores de potenciales.
  \item[\texttt{as(format, args...)}] Añade un mensaje descriptivo que
    aparece en caso de fallo. Muy útil en bucles.
  \item[\texttt{assertThatThrownBy(callable)}] Verifica que una lambda
    lanza una excepción del tipo esperado.
\end{description}

% ─────────────────────────────────────────────────────────────────────────────
\subsection{AssertJ Swing}
\label{sec:assertj-swing}

Librería para tests funcionales de interfaces Swing. Permite localizar
componentes, simular clicks y teclas, y verificar el estado visual. Se usa
en el módulo \texttt{org.openmarkov.gui} y en
\texttt{org.openmarkov.integrationtests}.

% ─────────────────────────────────────────────────────────────────────────────
\subsection{JaCoCo}
\label{sec:jacoco}

JaCoCo (\emph{Java Code Coverage}) instrumenta el bytecode para medir qué
líneas y ramas se ejecutan durante los tests. Se configura en el POM padre y
genera informes HTML en \texttt{target/site/jacoco/index.html} de cada módulo.

El script \texttt{test.sh} del workspace ofrece la opción \texttt{-r} para
abrir el informe de cobertura automáticamente tras ejecutar los tests.

% ─────────────────────────────────────────────────────────────────────────────
\subsection{Maven Surefire}
\label{sec:surefire}

Plugin de Maven que descubre y ejecuta los tests. Está configurado con la
versión 3.5.2 y usa el motor de JUnit Platform, que permite ejecutar
tanto tests de JUnit 5 como de jqwik en la misma pasada.


% ─────────────────────────────────────────────────────────────────────────────
\section{Herramientas incorporadas}
\label{sec:nuevas}

% ─────────────────────────────────────────────────────────────────────────────
\subsection{jqwik --- tests basados en propiedades}
\label{sec:jqwik}

jqwik\footnote{\url{https://jqwik.net/}} es un motor de \emph{property-based
testing} para Java que se integra con JUnit Platform. En lugar de verificar
un ejemplo concreto, una propiedad expresa un invariante que debe cumplirse
para \emph{cualquier} entrada válida. jqwik genera automáticamente cientos de
entradas aleatorias y, si encuentra un contraejemplo, lo reduce al caso
mínimo que reproduce el fallo (\emph{shrinking}).

\subsubsection{Anotaciones principales}

\begin{description}[leftmargin=2em, style=nextline]

  \item[\texttt{@Property}]
    Marca un método como test de propiedad. A diferencia de \texttt{@Test},
    que se ejecuta una sola vez, jqwik ejecuta un método \texttt{@Property}
    múltiples veces (por defecto 1000) con valores diferentes generados
    automáticamente para cada parámetro anotado con \texttt{@ForAll}.

    La idea fundamental es expresar un \textbf{invariante}: una condición
    que debe cumplirse para \emph{cualquier} entrada válida, no sólo para
    un ejemplo concreto. Si jqwik encuentra una combinación de valores que
    viola el invariante, reporta el contraejemplo y lo reduce automáticamente
    al caso mínimo que reproduce el fallo.

    En el siguiente ejemplo, el invariante es: <<el tamaño de la tabla de un
    potencial siempre es igual al producto del número de estados de todas sus
    variables>>. jqwik generará 1000 listas de variables con distintos números
    de estados y verificará que la propiedad se cumple en todos los casos.

\begin{lstlisting}[caption={\texttt{@Property}: invariante de tamaño de tabla}, label={lst:property}]
@Property
void tableSizeIsProductOfAllStateCounts(
        @ForAll("variableLists") List<Variable> variables) {
    TablePotential tp = new TablePotential(
            variables, PotentialRole.CONDITIONAL_PROBABILITY);
    int expected = variables.stream()
            .mapToInt(Variable::getNumStates)
            .reduce(1, (a, b) -> a * b);
    assertThat(tp.getTableSize()).isEqualTo(expected);
}
\end{lstlisting}

  \item[\texttt{@ForAll}]
    Se coloca delante de un parámetro de un método \texttt{@Property} para
    indicar a jqwik que debe generar un valor aleatorio para ese parámetro
    en cada ejecución.

    Para tipos primitivos (\texttt{int}, \texttt{double}, \texttt{String},
    etc.), jqwik sabe generar valores automáticamente. Se puede acotar el
    rango con anotaciones de restricción como \texttt{@IntRange} o
    \texttt{@DoubleRange}. Para tipos de dominio como \texttt{Variable} o
    \texttt{List<Variable>}, es necesario indicar un generador personalizado
    mediante \texttt{@ForAll("nombreDelGenerador")} (véase \texttt{@Provide}).

    En el siguiente ejemplo, jqwik genera automáticamente valores de
    \texttt{numStates} entre 1 y 50 sin necesidad de un generador
    personalizado:

\begin{lstlisting}[caption={\texttt{@ForAll} con tipos primitivos y restricciones}, label={lst:forall}]
@Property
void numStatesMatchesIntConstructor(
        @ForAll @IntRange(min = 1, max = 50) int numStates) {
    Variable v = new Variable("X", numStates);
    assertThat(v.getNumStates()).isEqualTo(numStates);
}
\end{lstlisting}

    Se pueden combinar múltiples parámetros \texttt{@ForAll} en una misma
    propiedad. jqwik genera todas las combinaciones:

\begin{lstlisting}[caption={\texttt{@ForAll} con múltiples parámetros}, label={lst:forall-multi}]
@Property
void reorderSwap_preservesValues(
        @ForAll @IntRange(min = 2, max = 5) int sA,
        @ForAll @IntRange(min = 2, max = 5) int sB) {
    Variable a = new Variable("A", sA);
    Variable b = new Variable("B", sB);
    // jqwik prueba con (sA=2,sB=2), (sA=3,sB=5), etc.
    // ...
}
\end{lstlisting}

  \item[\texttt{@Provide}]
    Marca un método que devuelve un \texttt{Arbitrary<T>} --- es decir,
    un generador de valores aleatorios de tipo \texttt{T} (véase la
    Sección~\ref{sec:generadores}). Este generador se conecta a un
    parámetro \texttt{@ForAll} por nombre: si el parámetro dice
    \texttt{@ForAll("variableLists")}, jqwik buscará un método anotado
    con \texttt{@Provide} que se llame \texttt{variableLists()}.

    IntelliJ marca estos métodos como no usados porque jqwik los invoca
    por reflexión; se recomienda añadir \texttt{@SuppressWarnings("unused")}.

    El siguiente generador produce listas de 1 a 4 variables, cada una con
    entre 2 y 5 estados, con nombres únicos (\texttt{V0}, \texttt{V1}, \ldots):

\begin{lstlisting}[caption={\texttt{@Provide}: generador de listas de variables}, label={lst:provide}]
@Provide
@SuppressWarnings("unused")
Arbitrary<List<Variable>> variableLists() {
    return Arbitraries.integers().between(1, 4).flatMap(numVars ->
        Arbitraries.integers().between(2, 5)
            .list().ofSize(numVars)
            .map(stateCounts -> {
                List<Variable> vars = new ArrayList<>(numVars);
                for (int i = 0; i < numVars; i++) {
                    vars.add(new Variable("V" + i,
                                          stateCounts.get(i)));
                }
                return vars;
            })
    );
}
\end{lstlisting}

    El generador se usaría así en una propiedad:

\begin{lstlisting}[caption={Uso de un generador personalizado en una propiedad}, label={lst:provide-usage}]
@Property
void valuesLengthEqualsTableSize(
        @ForAll("variableLists") List<Variable> variables) {
    TablePotential tp = new TablePotential(
            variables, PotentialRole.CONDITIONAL_PROBABILITY);
    assertThat(tp.getValues()).hasSize(tp.getTableSize());
}
\end{lstlisting}

  \item[\texttt{@Example}]
    Equivalente a \texttt{@Test} de JUnit, pero dentro de una clase
    gestionada por jqwik. Se usa para tests de regresión específicos
    que conviven con propiedades en la misma clase. A diferencia de
    \texttt{@Property}, un \texttt{@Example} se ejecuta exactamente una
    vez, con valores fijos elegidos por el desarrollador.

    Es especialmente útil para documentar bugs concretos que se han
    descubierto: el test de regresión asegura que el bug no reaparezca.
    En el siguiente ejemplo, se verifica que hacer undo de un
    \texttt{AddLinkEdit} no muta la lista de variables del potencial
    original de un nodo UTILITY:

\begin{lstlisting}[caption={\texttt{@Example}: test de regresión específico}, label={lst:example}]
@Example
void undoDoesNotMutateOldUtilityPotentialVariables()
        throws DoEditException {
    ProbNet net = new ProbNet(
            InfluenceDiagramType.getUniqueInstance());
    net.getPNESupport().setWithUndo(true);
    Variable vC = new Variable("C", 2);
    Variable vU = new Variable("U"); // utility
    new AddNodeEdit(net, vC, NodeType.CHANCE).executeEdit();
    new AddNodeEdit(net, vU, NodeType.UTILITY).executeEdit();

    AddLinkEdit edit = new AddLinkEdit(net, vC, vU, true);
    edit.executeEdit();
    edit.undo();

    List<Variable> varsAfterUndo =
            net.getNode(vU).getPotentials().getFirst()
               .getVariables();
    assertThat(varsAfterUndo).doesNotContain(vC);
}
\end{lstlisting}

\end{description}

\subsubsection{Anotaciones de restricción}

Estas anotaciones se combinan con \texttt{@ForAll} para acotar el rango de
los valores generados:

\begin{description}[leftmargin=2em, style=nextline]
  \item[\texttt{@IntRange(min, max)}]
    Restringe un \texttt{int} al rango \texttt{[min, max]}. Los valores
    por defecto son \texttt{min = 0}, \texttt{max = Integer.MAX\_VALUE},
    así que \texttt{@IntRange(max = 100)} es equivalente a
    \texttt{@IntRange(min = 0, max = 100)}.

  \item[\texttt{@DoubleRange(min, max)}]
    Restringe un \texttt{double} al rango especificado. Útil para factores
    de escala o probabilidades.
\end{description}

\subsubsection{API de generadores (\texttt{Arbitraries})}
\label{sec:generadores}

Un \textbf{generador} (tipo \texttt{Arbitrary<T>} en jqwik) es un objeto
que sabe producir valores aleatorios de tipo \texttt{T}. Cuando jqwik
ejecuta una propiedad, pide a cada generador que produzca un valor
distinto en cada iteración, usando un generador de números pseudoaleatorios
interno. Además, si el test falla, jqwik usa el generador para
\emph{reducir} (\emph{shrink}) el valor hacia el contraejemplo más
pequeño posible --- por ejemplo, reduciendo una lista de 5 elementos a
una de 2 si con eso basta para reproducir el fallo.

Para tipos primitivos como \texttt{int} o \texttt{String}, jqwik
proporciona generadores predefinidos. Para tipos de dominio como
\texttt{Variable} o \texttt{List<Variable>}, hay que construir generadores
personalizados componiendo los primitivos. La clase \texttt{Arbitraries}
es el punto de entrada para esta composición.

Los combinadores más usados en el proyecto son:

\begin{description}[leftmargin=2em, style=nextline]
  \item[\texttt{Arbitraries.integers().between(a, b)}]
    Genera enteros en el rango \texttt{[a, b]}, con distribución uniforme.
  \item[\texttt{Arbitraries.strings().alpha().ofMinLength(n).ofMaxLength(m)}]
    Genera cadenas compuestas sólo de letras, de longitud entre
    \texttt{n} y \texttt{m}.
  \item[\texttt{.flatMap(f)}]
    Encadena un generador con otro que \emph{depende} del valor generado.
    Imprescindible cuando un valor determina la estructura de otro --- por
    ejemplo, generar primero el número de variables y luego una lista de
    ese tamaño.
  \item[\texttt{.map(f)}]
    Transforma el valor generado sin cambiar la estructura del generador.
    Se usa para construir objetos de dominio (\texttt{Variable},
    \texttt{State[]}) a partir de valores primitivos.
  \item[\texttt{.list().ofSize(n)}]
    Genera una lista de exactamente \texttt{n} elementos del tipo base.
\end{description}

El Listado~\ref{lst:provide} muestra un ejemplo completo de composición:
\texttt{flatMap} genera primero el número de variables (1--4), luego
\texttt{.list().ofSize(numVars)} genera una lista de enteros de ese tamaño
(los números de estados), y finalmente \texttt{.map} transforma la lista
de enteros en una lista de objetos \texttt{Variable}.

\subsubsection{Fichero de semillas}

jqwik usa internamente un generador de números pseudoaleatorios para
producir los valores de cada ejecución. La \textbf{semilla}
(\emph{seed}) es el valor inicial de ese generador: dada la misma semilla,
jqwik produce exactamente la misma secuencia de valores, lo que permite
reproducir un fallo de forma determinista.

jqwik almacena las semillas de las ejecuciones anteriores en un fichero
\texttt{.jqwik-database} en la raíz de cada módulo. Si un test falló en
la ejecución anterior, jqwik reutiliza la misma semilla para verificar
que el fallo se ha corregido antes de probar semillas nuevas.

Este fichero es local a cada máquina de desarrollo y está excluido del
control de versiones mediante \texttt{.gitignore} en todos los módulos.

% ─────────────────────────────────────────────────────────────────────────────
\subsection{PITest --- testing de mutaciones}
\label{sec:pitest}

PITest\footnote{\url{https://pitest.org/}} es una herramienta de
\emph{mutation testing}. Introduce pequeñas modificaciones (\emph{mutantes})
en el código fuente compilado --- por ejemplo, cambiar un \texttt{+} por un
\texttt{-}, invertir una condición, o sustituir un valor de retorno --- y
ejecuta los tests. Si un test falla, el mutante ha sido \emph{eliminado}
(\emph{killed}); si todos los tests pasan, el mutante ha \emph{sobrevivido},
lo que indica que los tests no detectan esa alteración del código.

\subsubsection{Configuración}

PITest está configurado en el \texttt{pluginManagement} del POM padre y
activado con targets específicos en los módulos que lo usan (actualmente
\texttt{org.openmarkov.core}):

\begin{lstlisting}[caption={Configuración de PITest en el POM del módulo core},
  label={lst:pitest-config}, language=XML]
<plugin>
    <groupId>org.pitest</groupId>
    <artifactId>pitest-maven</artifactId>
    <configuration>
        <targetClasses>
            <param>org.openmarkov.core.*</param>
        </targetClasses>
        <targetTests>
            <param>org.openmarkov.core.*</param>
        </targetTests>
    </configuration>
</plugin>
\end{lstlisting}

\subsubsection{Ejecución e informe}

\begin{lstlisting}[caption={Ejecución de PITest}, label={lst:pitest-run},
  language=bash]
# Ejecutar mutation testing en el módulo core
mvn -f org.openmarkov.core/pom.xml pitest:mutationCoverage

# El informe HTML se genera en:
# org.openmarkov.core/target/pit-reports/YYYYMMDDHHMI/index.html
\end{lstlisting}

\subsubsection{Interpretación de resultados}

PITest genera un informe HTML que muestra, para cada clase:

\begin{itemize}
  \item El número de mutantes generados.
  \item Cuántos fueron eliminados (\emph{killed}) por los tests.
  \item Cuántos sobrevivieron (\emph{survived}) --- indican debilidades en
    los tests.
  \item La \emph{mutation score}: porcentaje de mutantes eliminados.
\end{itemize}

Un mutation score bajo no indica necesariamente que el código esté mal, sino
que los tests no son suficientemente sensibles a ciertos cambios. Es una
métrica complementaria a la cobertura de líneas de JaCoCo: un código puede
tener 100\% de cobertura de líneas pero un mutation score bajo si los tests
no verifican los valores correctos.

\subsubsection{Ejemplo ilustrativo}

Para entender cómo funciona PITest en la práctica, consideremos un método
sencillo que calcula el número de estados de una variable discreta:

\begin{lstlisting}[caption={Método original que será mutado por PITest},
  label={lst:pitest-original}]
public int getNumStates() {
    return states.size();
}
\end{lstlisting}

Un test que simplemente invoque el método sin verificar el resultado no
eliminará ningún mutante:

\begin{lstlisting}[caption={Test débil --- no elimina mutantes},
  label={lst:pitest-weak}]
@Test
void testGetNumStates() {
    Variable v = new Variable("X", 3);
    v.getNumStates(); // no hay asercion!
}
\end{lstlisting}

PITest genera varios mutantes a partir del código original. Por ejemplo,
el mutador \texttt{RETURN\_VALS} sustituye el valor de retorno por 0:

\begin{lstlisting}[caption={Mutante generado: valor de retorno reemplazado por 0},
  label={lst:pitest-mutant}]
// Mutante: replaced int return with 0
public int getNumStates() {
    return 0; // <-- mutacion introducida
}
\end{lstlisting}

El test débil del Listado~\ref{lst:pitest-weak} no detecta este cambio,
porque no comprueba el valor devuelto --- el mutante \emph{sobrevive}.
Un test que sí verifica el resultado elimina al mutante:

\begin{lstlisting}[caption={Test robusto --- elimina el mutante},
  label={lst:pitest-strong}]
@Test
void testGetNumStates() {
    Variable v = new Variable("X", 3);
    assertThat(v.getNumStates()).isEqualTo(3);
}
\end{lstlisting}

Con esta aserción, cuando PITest sustituye \texttt{states.size()} por
\texttt{0}, el test falla y el mutante queda eliminado (\emph{killed}).
Otros mutadores típicos de PITest incluyen:

\begin{itemize}
  \item \textbf{CONDITIONALS\_BOUNDARY}: cambia \texttt{<} por \texttt{<=}
    (y viceversa).
  \item \textbf{NEGATE\_CONDITIONALS}: invierte condiciones
    (\texttt{==} $\to$ \texttt{!=}).
  \item \textbf{MATH}: sustituye operadores aritméticos
    (\texttt{+} $\to$ \texttt{-}, \texttt{*} $\to$ \texttt{/}).
  \item \textbf{VOID\_METHOD\_CALLS}: elimina llamadas a métodos
    \texttt{void}.
  \item \textbf{EMPTY\_RETURNS}: devuelve colecciones vacías o cadenas vacías
    en lugar del valor real.
\end{itemize}

% ─────────────────────────────────────────────────────────────────────────────
\section{Convenciones del proyecto}
\label{sec:convenciones}

\begin{itemize}
  \item Las clases de test se nombran \texttt{*Test} para tests de ejemplo
    y \texttt{*PropertyTest} para tests basados en propiedades.
  \item Los tests de propiedades conviven con tests \texttt{@Example} en la
    misma clase cuando el ejemplo es una regresión del mismo componente.
  \item Los proveedores (\texttt{@Provide}) se agrupan al final de la clase
    en una sección separada por comentarios.
  \item Se usa \texttt{@SuppressWarnings("unused")} en los métodos
    \texttt{@Provide} para evitar advertencias de IntelliJ.
  \item Se declara \texttt{throws DoEditException} en las propiedades que
    ejecutan edits, en lugar de capturar la excepción.
  \item Los helpers de creación de redes (\texttt{freshBayesianNet()},
    \texttt{twoNodeBN()}) se definen como métodos privados o records
    al inicio de la clase de test.
\end{itemize}

% ─────────────────────────────────────────────────────────────────────────────
\section{Inventario de tests existentes}
\label{sec:inventario}

La Tabla~\ref{tab:tests} lista las clases de test basadas en propiedades
añadidas hasta la fecha, todas en el módulo \texttt{org.openmarkov.core}.

\begin{table}[H]
\centering
\begin{tabularx}{\textwidth}{lrX}
\toprule
\textbf{Clase de test} & \textbf{Props.} & \textbf{Qué verifica} \\
\midrule
\texttt{VariablePropertyTest} & 12 &
  Invariantes de nombre/baseName/timeSlice, índice de estados, constructor
  copia, potenciales delta \\
\texttt{TablePotentialPropertyTest} & 11 &
  Tamaño de tabla, offsets, \texttt{scalePotential}, constructor copia,
  \texttt{setUniform} \\
\texttt{TablePotentialReorderProject\-PropertyTest} & 10 &
  \texttt{reorder()} preserva valores para 2 y 3 variables,
  identidad, round-trip; \texttt{tableProject()} con evidencia parcial,
  total y nula \\
\texttt{AddNodeEditPropertyTest} & 9 &
  Ciclo do/undo/redo de \texttt{AddNodeEdit}, rechazo de nombres
  duplicados \\
\texttt{AddLinkEditPropertyTest} & 11 &
  Ciclo do/undo/redo de \texttt{AddLinkEdit}, rechazo de
  self-loop y links duplicados, regresión de mutación de potencial
  UTILITY \\
\texttt{ProbNetOperationsBarren\-NodesPropertyTest} & 12 &
  \texttt{removeBarrenNodes()}: seguridad de nodos de interés/evidencia,
  eliminación de hojas/cadenas/diamantes, idempotencia,
  completitud \\
\midrule
\textbf{Total} & \textbf{65} & \\
\bottomrule
\end{tabularx}
\caption{Clases de test basadas en propiedades (jqwik)}
\label{tab:tests}
\end{table}

\end{document}
